\documentclass[11pt]{article}
\usepackage[margin=1in]{geometry}
\usepackage{amsmath,amssymb,amsthm}
\usepackage{booktabs}
\usepackage{float}
\usepackage{tabularx,array}
\newcolumntype{Y}{>{\raggedright\arraybackslash}X}
\usepackage[hyphens]{url}
\usepackage{hyperref}
\setlength{\parskip}{3pt}
\setlength{\emergencystretch}{3em}\sloppy  % long code names: stretch the spaces instead of overflowing

% Jexpresso notation
\newcommand{\ngl}{\mathtt{ngl}}
\newcommand{\nelem}{\mathtt{nelem}}
\newcommand{\npoin}{\mathtt{npoin}}
\newcommand{\iel}{\mathtt{iel}}
\newcommand{\connijk}{\mathtt{connijk}}
\newcommand{\Je}{J^{e}}
\newcommand{\K}{\mathbf{K}}
\newcommand{\M}{\mathbf{M}}
\newcommand{\uu}{\mathbf{u}}
\newcommand{\bb}{\mathbf{b}}
\newcommand{\rr}{\mathbf{r}}
\newcommand{\zz}{\mathbf{z}}
\newcommand{\Khat}{\widehat{K}}
\newcommand{\What}{\widehat{\omega}}
\newcommand{\code}[1]{\texttt{#1}}
\newtheorem{remark}{Remark}

\title{\vspace{-2.5em}Iterative solvers for the 3D periodic SEM Poisson problem in Jexpresso:\\
SEM AMG, SC AMG, p-MG AMG and p-MG GMG}
\author{Jexpresso --- \code{sm/elementLearning}}
\date{}

\begin{document}
\maketitle
\vspace{-2.5em}

\section{The discrete problem}\label{sec:problem}

The deck \code{problems/Elliptic/poisson\_periodic\_sem\_3d} solves $-\nabla^2 u = f$ on the
periodic box $\Omega=[0,2\pi]^3$ with \code{sem\_setup}. The mesh has $\nelem = n_e^3$ hexahedra,
each with $\ngl = N+1$ Legendre--Gauss--Lobatto (LGL) points $\xi_i$ per direction, with weights $\omega_i$.
The basis is the tensor product of the 1D Lagrange polynomials $\psi_i(\xi)$, whose
derivatives $\psi'_i(\xi_k)$ are Jexpresso's \code{basis.dpsi[i,k]}. Local node $(i,j,k)$ of
element $\iel$ has global number $\connijk[\iel,i,j,k]$.

The quadrature is collocated (\code{Inexact}, $Q=N$), so the mass matrix is diagonal:
\begin{equation}
  \M_{I} \;=\; \sum_{\iel}\ \sum_{(i,j,k)\mapsto I} \omega_i\,\omega_j\,\omega_k\,\Je_{ijk},
  \label{eq:mass}
\end{equation}
where $\Je_{ijk}$ is the Jacobian of element $\iel$ at node $(i,j,k)$ (Jexpresso's \code{metrics.Je}). The
stiffness matrix is the direct stiffness summation (DSS) of the element matrices
(\code{DSS\_laplace\_sparse\_3D}):
\begin{equation}
  \K_{IJ} \;=\; \sum_{\iel}\ \sum_{l,m,n} \omega_l\omega_m\omega_n\,\Je_{lmn}\;
  \nabla\psi_I(\xi_l,\eta_m,\zeta_n)\cdot\nabla\psi_J(\xi_l,\eta_m,\zeta_n).
  \label{eq:stiff}
\end{equation}
Here $\nabla = (\xi_x\partial_\xi + \eta_x\partial_\eta + \zeta_x\partial_\zeta,\ \dots)$, with the
inverse metric terms \code{d$\xi$dx}, \dots, \code{d$\zeta$dz}. Periodicity identifies the nodes of
opposite faces into classes. The solve uses the class numbering, $\K = P^{T} L P$
(\code{periodic\_sem\_system}), and $\K$ is singular, with the constants in its null space.
The right-hand side $\bb = \M f$ is compatible ($\mathbf{1}^T\bb = 0$), and the solution is
fixed by a zero $\M$-weighted mean.

\paragraph{Tensor structure on box elements.} On an axis-aligned box element the metric terms are
constant, $\xi_y=\xi_z=\eta_x=\dots=0$. Equation~\eqref{eq:stiff} then factors into 1D matrices:
\begin{equation}
  A_{\iel} \;=\; c^{\iel}_x\, \What\otimes\What\otimes\Khat \;+\; c^{\iel}_y\, \What\otimes\Khat\otimes\What
  \;+\; c^{\iel}_z\, \Khat\otimes\What\otimes\What ,
  \label{eq:tensor}
\end{equation}
(the index $i$ runs fastest). Here
\begin{equation}
  \Khat_{ab} = \sum_k \omega_k\,\psi'_a(\xi_k)\,\psi'_b(\xi_k), \qquad \What=\mathrm{diag}(\omega_k),
  \qquad c^{\iel}_x = \Je\,\xi_x^2,\quad c^{\iel}_y = \Je\,\eta_y^2,\quad c^{\iel}_z = \Je\,\zeta_z^2 .
  \label{eq:1d}
\end{equation}
Because $\What$ is diagonal, each row of $A_{\iel}$ couples a node only to the nodes on the three grid
lines through it. After assembly a row of $\K$ has between $3N+1$ nonzeros (a node inside an element)
and $6N+1$ (an element vertex, whose three lines each cross two elements). Two of the four solvers below (SC AMG and the two p-MG solvers) exploit~\eqref{eq:tensor}.

All four solvers run preconditioned conjugate gradients (CG) on the same matrix $\K$. They stop at
a relative preconditioned residual $\sqrt{\rr^T\zz}\,/\sqrt{\rr_0^T\zz_0}\le 10^{-12}$ and agree with the
sparse direct solve to round-off (\code{verify\_3d\_jexpresso.jl}). They differ only in the
preconditioner, and in whether CG works on all of $\K$ or only on its skeleton.

\section{The four solvers}

\subsection{SEM AMG (\code{sem\_amg}, \code{:linsolve\_amg => true})}
\begin{itemize}
  \item \textbf{Operator:} the full system $\K\uu=\bb$, with one class pinned: CG on $\K_{2:,2:}$ (\code{periodic\_sem\_amg\_solve}).
  \item \textbf{Preconditioner:} one V-cycle of smoothed-aggregation AMG (AlgebraicMultigrid.jl,
        \code{jx\_amg\_setup}), built from the entries of $\K$ alone.
  \item \textbf{Cost:} the matrix--vector product in CG is threaded (\code{JXSymCSC}); the AMG setup and
        V-cycle are serial.
\end{itemize}
AMG knows nothing about the SEM. Aggregation-based AMG is designed for M-matrices with
connectivity like a low-order stencil. The high-order matrix $\K$ is different:
\begin{itemize}
  \item the GLL stiffness $\Khat$ has positive off-diagonal entries;
  \item the strongest couplings run along the GLL lines;
  \item the nodes cluster near element boundaries, $\Delta\xi\sim N^{-2}$.
\end{itemize}
The aggregates therefore approximate the smooth modes poorly, and the iteration count grows with
$N$: 29 at $4^3$ elements, $N=4$; 63 at $8^3$, $N=6$; 59 at $16^3$, $N=4$.

\subsection{SC AMG (\code{sc\_amg}, \code{:lstatic\_condensation => true}, \code{:EL\_skeleton\_solver => "amg"})}
Split the unknowns into element interiors $o$ ($(N-1)^3$ per element) and the skeleton $b$: the
element-boundary nodes, $\partial O$ in the element-learning notation. Eliminating the interiors
gives the Schur complement on the skeleton:
\begin{equation}
  S\,\uu_b = \bb_b - \sum_{\iel} A_{bo}A_{oo}^{-1}\bb_o,\qquad
  S = A_{bb} - \sum_{\iel} A_{bo}\,A_{oo}^{-1}A_{ob},\qquad
  \uu_o = A_{oo}^{-1}(\bb_o - A_{ob}\uu_b).
  \label{eq:schur}
\end{equation}
Here $A_{oo}$ is the interior block of element $\iel$ (Jexpresso's \code{Avovo}), and $A_{ob}$ the
interior--boundary block (\code{Avovb}). The default implementation, \code{el\_sc\_schur\_cg!},
never forms $S$:
\begin{itemize}
  \item \textbf{Interior solves by fast diagonalization.} Take the interior part of~\eqref{eq:tensor},
        $A_{oo} = c_x\What_o\otimes\What_o\otimes\Khat_o + \dots$ (rows and columns $2{:}N$ of the 1D matrices).
        Solve the 1D generalized eigenproblem $\Khat_o s = \lambda\What_o s$, with $S_1^T\What_o S_1=I$. Then
        \begin{equation}
          A_{oo}^{-1} = (S_1\otimes S_1\otimes S_1)\,\mathrm{diag}\!\Big(\frac{1}{c_x\lambda_i+c_y\lambda_j+c_z\lambda_k}\Big)\,(S_1\otimes S_1\otimes S_1)^T,
          \label{eq:fdm}
        \end{equation}
        which costs $O(N^4)$ per element. Only $(c_x,c_y,c_z)$ is stored per element; these are read off
        $\K$ and checked entry by entry against~\eqref{eq:tensor}.
  \item \textbf{Matrix-free Schur complement:} $S\mathbf{p} = \big(\K\,[\mathbf{p};\,-A_{oo}^{-1}A_{ob}\mathbf{p}]\big)_b$,
        one sparse matvec plus the interior solves~\eqref{eq:fdm}.
  \item \textbf{Preconditioner:} the SEM AMG V-cycle of the \emph{full} $\K$, $P^{-1}$, restricted to the skeleton:
        $\zz_b = (P^{-1}[\rr_b;0])_b$. For an SPD matrix, $(\K^{-1})_{bb}=S^{-1}$ exactly. So if
        $c_1\K\le P\le c_2\K$, then $(P^{-1})_{bb}$ is spectrally equivalent to $S^{-1}$ with the same
        constants, and $\kappa\big((P^{-1})_{bb}S\big)\le\kappa(P^{-1}\K)$.
\end{itemize}
SC AMG therefore inherits SEM AMG's convergence: 57 iterations against 59--63, a few fewer because
CG runs on the smaller space. Each iteration costs one matvec, the interior solves and one full
V-cycle. Its setup is the same AMG hierarchy as SEM AMG's.

The alternative, \code{:EL\_sc\_amg => "assembled"}, assembles $S$ (\code{el\_sc\_solve!}) and gives it to AMG. That
needs fewer iterations (34 at $8^3$, $N=6$), but $S$ is far denser than $\K$. A face node of the skeleton couples
to all the active boundary nodes of both neighbouring elements, about 10 times the nonzeros of $\K$ at
$N=8$. Each iteration, and the setup, therefore cost much more.

\subsection{p-MG + AMG and p-MG + GMG (\code{pmg\_amg}, \code{pmg\_gmg}, \code{:linsolve\_pmg => "amg" | "gmg"})}
These run CG on the full singular $\K$, preconditioned by one p-multigrid V-cycle (\code{pmultigrid.jl},
\code{jx\_pmg\_setup}, \code{jx\_pmg\_cg}). The hierarchy follows the SEM order rather than the matrix entries.

\paragraph{p-levels.} The orders are $p_0=N,\ p_{\ell+1}=\lfloor p_\ell/2\rfloor$, down to $p=1$ (for example $8\to4\to2\to1$), all on
the same $\nelem$ elements:
\begin{itemize}
  \item $\ell=0$: $\K$ itself, built by \code{sem\_setup}.
  \item $\ell>0$: a rediscretization of the same operator with the LGL basis of order $p$, built by Jexpresso:
        \code{basis\_structs\_$\xi$\_$\omega$!(LGL(), p)} gives $\xi^{(p)}_k,\omega^{(p)}_k$;
        \code{build\_Interpolation\_basis!} gives $\psi^{(p)\prime}_a(\xi^{(p)}_k)$; the coefficients $c^{\iel}$
        come from \code{metrics}. $A^{(p)}$ is the DSS of~\eqref{eq:tensor} with $\Khat^{(p)},\What^{(p)}$.
  \item \textbf{Check:} the order-$N$ form of~\eqref{eq:tensor} must reproduce $\mathrm{diag}(\K)$, and each element's
        local axes must map to $\pm x,\pm y,\pm z$.
\end{itemize}

\paragraph{Transfers.} Prolongation from order $p_c$ to $p_f$ is the tensor-product Lagrange interpolation
\begin{equation}
  \Pi_{kl} = \psi^{(p_c)}_{l}\big(\xi^{(p_f)}_{k}\big),\qquad
  \uu^{(p_f)}_{\iel} = (\Pi\otimes\Pi\otimes\Pi)\,\uu^{(p_c)}_{\iel},
  \label{eq:transfer}
\end{equation}
applied element by element in three 1D passes. Every fine node is written by one owning element. Restriction is
the transpose. It accumulates element by element, with the elements coloured so that the threads never
write the same coarse node.

\paragraph{Smoother.} On every level, Chebyshev--Jacobi of degree 3 targets the upper part
$[0.25\,\lambda_{\max},\,1.1\,\lambda_{\max}]$ of the spectrum of $D^{-1}A^{(p)}$. Here $D=\mathrm{diag}(A^{(p)})$,
and $\lambda_{\max}$ comes from 20 power iterations. Each Chebyshev step is
\[
  \mathbf{x}\leftarrow\mathbf{x}+\mathbf{d},\qquad
  \mathbf{r}\leftarrow\mathbf{r}-D^{-1}A\mathbf{d},\qquad
  \mathbf{d}\leftarrow\rho_{k+1}\rho_k\,\mathbf{d}+\tfrac{2\rho_{k+1}}{\delta}\,\mathbf{r}.
\]
It needs only threaded matvecs, and it is a fixed polynomial in $D^{-1}A$, so pre- and post-smoothing keep
the V-cycle symmetric, as CG requires.

\paragraph{Why this works for high order.} A Jacobi-type smoother removes the components of the error
that oscillate on the scale of the GLL spacing. What remains is well represented by a polynomial of
half the order: the same operator, rediscretized with $p/2$, is spectrally close on that subspace.
The hierarchy ends at $p=1$: the trilinear SEM with two-point GLL (trapezoidal) quadrature. There,
$\Khat^{(1)}=\tfrac12\begin{bmatrix}1&-1\\-1&1\end{bmatrix}$ and $\What^{(1)}=I$, so~\eqref{eq:tensor} is the
7-point Laplacian on the element vertices: a symmetric M-matrix on an $n_e^3$ grid, exactly the problem
AMG and geometric multigrid are designed for.

\paragraph{The $p=1$ level: the only difference between the two variants.}
\begin{itemize}
  \item \textbf{p-MG + AMG:} one smoothed-aggregation AMG V-cycle on $A^{(1)}$, with one node pinned.
  \item \textbf{p-MG + GMG:} geometric h-multigrid on the periodic $n_e^3$ vertex grid:
        \begin{itemize}
          \item trilinear interpolation $P_h$ (weights $1$ and $\tfrac12$ per direction);
          \item Galerkin coarse operators $A_{2h}=P_h^TA_hP_h$;
          \item coarsening by 2 while every dimension is even and the grid has more than 512 nodes;
          \item a sparse Cholesky solve of the pinned coarsest operator.
        \end{itemize}
\end{itemize}
Both variants share the p-levels, so the iteration counts are nearly equal (12 against 10 at $16^3$, $N=4$). GMG needs
no aggregation setup and its transfers are known in advance, so its setup is cheaper. AMG needs no
structured grid.

\paragraph{The singular system.} CG runs on $\K$ itself. The V-cycle subtracts the mean from its input and output.
Interpolation reproduces constants, so coarse right-hand sides stay compatible. The coarsest solves pin one node,
which solves a compatible singular system exactly.

\subsection{Side by side}

\begin{table}[H]
\centering\small
\begin{tabularx}{\linewidth}{@{}lYYYY@{}}
\toprule
 & SEM AMG & SC AMG & p-MG + AMG & p-MG + GMG \\
\midrule
CG operator & $\K$ (pinned) & $S$ on the skeleton, matrix-free & $\K$ (singular) & $\K$ (singular)\\
preconditioner & AMG V-cycle of $\K$ & AMG V-cycle of $\K$, restricted to $b$ & p-V-cycle $N\to\dots\to1$, AMG at $p=1$ & p-V-cycle $N\to\dots\to1$, GMG at $p=1$ \\
uses the SEM structure & no & yes: $A_{oo}^{-1}$ by fast diagonalization~\eqref{eq:fdm} & yes: levels~\eqref{eq:tensor}, transfers~\eqref{eq:transfer} & yes, and the vertex grid\\
smoother & Gauss--Seidel (AMG) & Gauss--Seidel (AMG) & Chebyshev--Jacobi; AMG at $p=1$ & Chebyshev--Jacobi on all levels\\
threaded & matvec & matvec, interior solves & all p-levels & all levels\\
needs & any SPD $\K$ & box elements, else assembled $S$ & box elements, periodic box & same, plus $n_e$ divisible by 2\\
\bottomrule
\end{tabularx}
\caption{The four iterative solvers. ``Box elements'' means axis-aligned hexahedra with constant metric terms, the
structure assumed by~\eqref{eq:tensor}.}
\end{table}

\begin{table}[H]
\centering\small
\begin{tabular}{@{}lrrrrrrrr@{}}
\toprule
 & \multicolumn{4}{c}{$8^3$ elements, $N=6$ ($1.1\cdot10^5$ unknowns)} & \multicolumn{4}{c}{$16^3$ elements, $N=4$ ($2.6\cdot10^5$ unknowns)}\\
\cmidrule(lr){2-5}\cmidrule(l){6-9}
 & iters & setup [s] & solve [s] & & iters & setup [s] & solve [s] & \\
\midrule
SEM AMG     & 63 & 3.1 & 3.5 & & 59 & 7.6 & 5.3 & \\
SC AMG      & 57 & 2.6 & 3.2 & & 57 & 8.8 & 5.9 & \\
p-MG + AMG  & 13 & 1.3 & 0.8 & & 12 & 3.2 & 1.1 & \\
p-MG + GMG  & 13 & --  & --  & & 10 & 1.3 & 0.8 & \\
\bottomrule
\end{tabular}
\caption{Local measurements, 4 threads on a shared 4-core container; the timings are noisy, the iteration
counts are not. ``setup'' includes the periodic reduction (0.3--2\,s). The p-MG + GMG timings at $8^3$
were not recorded. The Wulver sweep (\code{run\_wulver3d.sbatch}) gives the definitive numbers.}
\end{table}

\section{Why there is no SC GMG}\label{sec:noscgmg}

\paragraph{1. Geometric multigrid needs a grid hierarchy for the operator it solves.}
GMG rests on three geometric ingredients:
\begin{itemize}
  \item a sequence of meshes on which the \emph{same} discrete operator can be rebuilt;
  \item transfers that interpolate between them;
  \item a local smoother that removes the high-frequency error on each.
\end{itemize}
The p-MG levels and the $p=1$ vertex grid have all three, because every level discretizes $-\nabla^2$
with the SEM. The Schur complement $S$ in~\eqref{eq:schur} has none of them:
\begin{itemize}
  \item It lives only on the skeleton, the union of element faces, which is not a volume mesh.
  \item It is the discrete Steklov--Poincar\'e (Dirichlet-to-Neumann) operator of the elements: dense within each
        element's boundary and nonlocal, with a condition number growing like $h^{-1}$ rather than $h^{-2}$.
  \item A coarser skeleton is not the skeleton of a coarser SEM. Halving $n_e$ turns half of the fine skeleton
        into interiors of the coarse elements, so the coarse Schur complement eliminates nodes that are unknowns
        on the fine level.
  \item The transfers between two skeleton levels must therefore include the discrete harmonic extension
        $-A_{oo}^{-1}A_{ob}$ into the eliminated interiors. That is substructuring (balancing domain
        decomposition, BDDC, Schur-complement multigrid), not geometric multigrid.
  \item A pointwise smoother on $S$ needs $\mathrm{diag}(S)$. That is only available element by element
        through $A_{bo}A_{oo}^{-1}A_{ob}$, the work the matrix-free path avoids.
\end{itemize}
A geometric multigrid applied to $S$ itself is therefore a different method, not a variant of the ones implemented.

\paragraph{2. Its benefit is already available through the restriction identity.}
SC AMG never applies AMG to $S$. It restricts the full-system V-cycle, $\zz_b=(P^{-1}[\rr_b;0])_b$, which is a
valid preconditioner for $S$ \emph{whatever} multigrid $P$ is, because $(\K^{-1})_{bb}=S^{-1}$. An ``SC GMG'' in
the same spirit would replace $P$ by the p-MG + GMG V-cycle:
\begin{itemize}
  \item it would inherit p-MG's iteration counts, a few fewer than on the full system;
  \item each iteration would cost one p-MG V-cycle (on the full grid) plus one sparse matvec and the interior solves~\eqref{eq:fdm} for $S\mathbf{p}$;
  \item p-MG + GMG on $\K$ costs one V-cycle plus one matvec.
\end{itemize}
On this problem it would solve the same system with the same preconditioner and a little more work per
iteration, so it cannot beat \code{pmg\_gmg}. Condensation pays off in the SC \emph{direct} solver, where it
removes the interior fill from the factorization, and in element learning, where the per-element operators
$T^{\iel}=A_{oo}^{-1}A_{ob}$ are the quantity predicted by the network. It does not pay off with a multigrid
preconditioner that already handles the interior modes through its smoother.

\paragraph{3. In code.} This preconditioner can be added: \code{\_sc\_pcg} calls the preconditioner only through
\code{ldiv!}, and \code{JXPMG} implements \code{ldiv!} on the full singular system. It was not added, so that
\code{sc\_amg} stays comparable with the earlier runs, and because of point~2.

\end{document}
